% 使用大连理工大学学位论文模板类
% twoside = 双面排版
% hide    = 隐藏某些信息，通常用于盲审/匿名版
\documentclass[twoside,nohide]{DUT-thesis-grd}

%============== 更改数学字体设置（当前默认未启用） ==============
% 下面两行现在是注释状态，不会生效
% 如果你觉得默认数学字体太细，可以取消注释试试
% unicode-math：支持 Unicode 数学输入，也能更方便切换数学字体
% XITS Math：一种较常见、较适合论文的数学字体

% \usepackage[bold-style=ISO]{unicode-math} % 使用 unicode-math 宏包
% \setmathfont{XITS Math}                    % 设置数学字体为 XITS Math

%================ 用户信息填写（value 宏） ================
% 封面与学位信息
\title{论文题目}
\author{姓名}
\studentnumber{学号}
\advisor{导师}
\practiceadvisor{实践导师}
\major{专业}
\defenddate{2026年05月20日}
% 声明页签名图片（作者签名、导师签名）
\authorsignature{figures/figure1}
\supervisorsignature{figures/figure2}

% 答辩委员会与答辩信息
\committeeTitleOne{大连理工大学教授:}
\committeeTitleTwo{大连理工大学副教授:}
\committeeTitleThree{XXXXXXXXXXXX高工:}
\committeeTitleFour{XXXXXXXX大学教授:}
\committeeTitleFive{大连理工大学研究员:}

\committeeMemberOne{XXX}
\committeeMemberTwo{XXX}
\committeeMemberThree{XXX}
\committeeMemberFour{XXX}
\committeeMemberFive{XXX}

\defenseLocation{大连理工大学开发区校区}

%================ 进入正文 ================
\begin{document}

% 生成封面页
\maketitle

% 生成论文原创性声明和使用授权页
\makeDeclareOriginal

% 生成答辩委员会信息页
\committee

%================ 进入前置部分 ================
\frontmatter

% 插入摘要文件
\include{chapters/abstract}

% 生成中文目录
\tableofcontents

% 生成英文目录
\tableofengcontents

% 临时去掉图表目录里“各章之间的额外垂直间距”，让目录更紧凑
\let\origaddvspace\addvspace
\renewcommand{\addvspace}[1]{}

% 把“图目录”加入总目录
\addcontentsline{toc}{chapter}{图目录}
% 生成图目录
\listoffigures

% 把“表目录”加入总目录
\addcontentsline{toc}{chapter}{表目录}
% 生成表目录
\listoftables

% 恢复 addvspace 的原始定义
\renewcommand{\addvspace}[1]{\origaddvspace{#1}}

% 插入符号对照表
\input{chapters/denotation}

%================ 进入正文部分 ================
\mainmatter

% 临时关闭“空白页”，这样封面以及目录等会插入空白页而章节之间不会自动插入空白页
\DUTDisableBlankPage

% 按插入各章正文
\include{chapters/chapter1}
\include{chapters/chapter2}
\include{chapters/chapter3}
\include{chapters/chapter4}
\include{chapters/chapter5}
\include{chapters/chapter6}

% 参考文献
% 如果你有多个 .bib 文件，可以用下面这种写法（当前被注释）
%\bibliography{reference/chap1,reference/chap2}
\bibliography{reference/chap1}

%% 附录
%% appendx 之后，章节会重新编号，通常变成 A、B、C...
\appendix

% 把附录章节名改成“附录 A / 附录 B ...”
\renewcommand{\appendixname}{附录~\Alph{chapter}}

% 设置 ctex 的 chapter 名称为“附录”
\ctexset{chapter={name={附录}}}

% 附录中的图、表、公式、代码编号改成 A.1 / A.2 这种形式
\renewcommand{\thefigure}{\Alph{chapter}.\arabic{figure}}
\renewcommand{\thetable}{\Alph{chapter}.\arabic{table}}
\renewcommand{\theequation}{\Alph{chapter}.\arabic{equation}}
\renewcommand{\thelstlisting}{\Alph{chapter}.\arabic{lstlisting}}

% 设置附录中表格标题名称
\renewcommand{\tablename}{附录-表}
% 设置英文第二标题中的表名（模板双语 caption 用）
\captionsetup[table][bi-second]{name=App.Tab.}

% 设置附录中图片标题名称
\renewcommand{\figurename}{附录-图}
% 设置英文第二标题中的图名
\captionsetup[figure][bi-second]{name=App.Fig.}

% 插入附录 A、附录 B
\include{chapters/app1} 
\include{chapters/app2} 

%================ 进入后置部分 ================
% 一般这后面的章节不再按正文方式编号
\backmatter

% 插入发表成果部分
\include{chapters/pub}

% 插入致谢部分
\include{chapters/thanks}

% 文档结束
\end{document}
